\section{Energy in the magnetic field}
\begin{frame} \ftx{Energy in the magnetic field} 
	\s{To calculate the energy stored in a coil, the toroidal coil illustrated in Figure \ref{AbbSpule1} is connected to a DC voltage source. At time $t = 0$, the voltage source is switched on. The voltage across the coil is therefore equal to the source voltage $U$ at every point in time $t > 0$. 
    The current flowing through the coil will increase linearly according to equation \ref{GLInduktivitaet2}. The energy supplied in this way generates a flux density in the coil core. 
 	To calculate the energy supplied to the coil, the known equation for electrical power \ref{GlArbeit} can be transformed and then the voltage can be replaced by equation \ref{GLInduktivitaet2}.
    As a result of this transformation, the change in magnetic energy $\mathrm{d}W_{\mathrm{L}}$ is obtained, which is calculated from the multiplication of inductance $L$, the current through the inductance $i_{\mathrm{L}}$ and the rate of change of the current through the inductance $\mathrm{d}i_{\mathrm{L}}$ (see equation \ref {GlÄnderungmagnetischeEnergie}).
	
\renewcommand{\figurename}{Figure}		
	% \f{width=0.5\textwidth}{width=0.5\textwidth}{Spule}{Ringkernspule mit einer angelegten Spannung \label{AbbSpule1}{\tiny(Quelle: \cite[Abbildung 6.26]{Albach})}}	

		\fo{width=0.35\textwidth}{width=0.35\textwidth}{Ringspule_MagFeld}{
			\put(6, 53){\color{current}$I$}
			\put(11, 1){\color{current}$I$}
			\put(76, 15){\color{magnetfeld}$\ell_{\mathrm{m}}$}
		}{{\bf Closed toroidal coil.} \label{AbbSpule1}}	
	}	
	\b{
		\begin{figure}[htbp]
			\centering
			\begin{minipage}{0.4\textwidth}
				\centering
				\fo{width=0.8\textwidth}{width=0.8\textwidth}{Ringspule_MagFeld}{
					\put(6, 53){\color{current}$I$}
					\put(11, 1){\color{current}$I$}
					\put(76, 15){\color{magnetfeld}$\ell_{\mathrm{m}}$}
				}{\bf Geschlossene Ringkernspule.}
			\end{minipage}
			\hfill
			\begin{minipage}{0.54\textwidth}
				\begin{eqa}
					\onslide<1->{p &= \frac{\d W}{\d t} = u\cdot i \qquad u = L\cdot\frac{\d i}{\d t}}\\
					\onslide<2->{\d W_{\mathrm{m}} &= u_{\mathrm{L}}\cdot i_{\mathrm{L}} \cdot \d t = L\cdot \frac{\d i}{\d t}\cdot i_{\mathrm{L}}\cdot \d t = L\cdot i_{\mathrm{L}} \cdot \d i_{\mathrm{L}}}
				\end{eqa}
				
				\begin{eqa}
				\begin{aligned}
					\onslide<3->{W_{\mathrm{m}} &= L\cdot \int\limits_0^I i_{\mathrm{L}}\cdot \mathrm{d}i_{\mathrm{L}} = \frac{1}{2}\cdot L\cdot I^2 \\
					&= \frac{1}{2}\cdot N \cdot \varPhi \cdot I}
				\end{aligned}
				\end{eqa}
			\end{minipage}
		\end{figure}
	}
	\s{
		\begin{eqa}				
				p &= \frac{\d W}{\d t} = u\cdot i \label{GlArbeit}\\
				u &= L\cdot\frac{\d i}{\d t} \tag{\ref{GLInduktivitaet2}}\\
				\d W_{\mathrm{m}}&=u_{\mathrm{L}}\cdot i_{\mathrm{L}} \cdot \d t = L\cdot \frac{\d i}{\d t}\cdot i_{\mathrm{L}}\cdot \d t = L\cdot i_{\mathrm{L}} \cdot \d i_{\mathrm{L}} \label{GlÄnderungmagnetischeEnergie}
		\end{eqa}

	The total energy is the integral over $\d W_{\mathrm{m}}$ from the initial value $i_{\mathrm{L}}=0$ to the static final value $i_{\mathrm{L}} = I$.
	Assuming a constant inductance $L$, the following applies:
	
	\begin{equation}
	W_m = L \cdot \int_0^{I_L} i_L \cdot \d i_L = \frac{1}{2} \cdot L \cdot I^2 = \frac{1}{2} \cdot N \cdot \varPhi \cdot I
	\end{equation}}
	\nomenclature[FF]{$F$}{Kraft \nomunit{N}}
	\nomenclature[FP]{$P$}{Wirkleistung \nomunit{W}}
	\nomenclature[FW]{$W$}{Energie \nomunit{J}}

\end{frame}
\begin{frame} \ftx{Energy in the magnetic field} 
	\s{The magnetic energy can also be calculated from the field magnitudes by integration with the magnetic flux density:
	\begin{equation}
		W_{\mathrm{m}} = \ell_{\mathrm{m}}\cdot A\cdot \int\limits_0^{B_L} H\d B = \ell_{\mathrm{m}}\cdot A \cdot \frac{B_{\mathrm{L}}^2}{2\cdot \mu_{\mathrm{r}}\cdot \mu_0} 
	\end{equation}}
	\b{Calculation based on field sizes:
	\begin{equation}
		\begin{aligned}
		\onslide<1->{W_{\mathrm{m}} &= \ell_{\mathrm{m}}\cdot A\cdot \int\limits_0^{B_L} H\d B \\
		  &= \ell_{\mathrm{m}}\cdot A \cdot \frac{B_{\mathrm{L}}^2}{2\cdot \mu_{\mathrm{r}}\cdot \mu_0}} \\[10pt]
		\onslide<2->{\frac{W_{\mathrm{m}}}{V} &= \frac{B_{\mathrm{L}}^2}{2\cdot \mu_0} = \frac{F}{A}}
		\end{aligned}
	\end{equation}}
	\pause
	\s{In an electromagnet, the energy density in the air gap (i.e. the energy per volume of the air gap) is particularly important: it represents the \glqq force of the electromagnet\grqq. The maximum force is exerted precisely at the transition from the air gap to the iron core – corresponding to a theoretically infinitely narrow air gap. Since the air gap, as the name suggests, is usually filled with air, $\mu_{\mathrm{r}} = 1$ in normal cases and can be omitted.
	
	\begin{eq}
		\frac{W_{\mathrm{m}}}{V} = \frac{B_{\mathrm{L}}^2}{2\cdot \mu_0} = \frac{F}{A}
	\end{eq}
	}
\end{frame}